\scanpage{22}

\uuline{NB} $1^\circ$. La corr. $(E,\Delta,Q)\longmapsto\deleted{\{}\mathcal A_Q\deleted{\unsure{\mathcal A_Q\subset|\mathcal A_Q\to}\underline{\operatorname{End}}(E)}$ est \uline{fonctorielle}\inserted{(pour les iso)}, et l'hom (pas néc. inj.)

\[
\mathcal A_Q\longrightarrow\underline{\operatorname{End}}(E)
\]

\deleted{est} fonctoriel. De plus, ces constructions commu\deleted{\unsure{tent}} aux changements de base.

$2^\circ$. Il en résulte, par transport de structure, que \inserted{pour $\lambda\in\Gamma\,\underline O_X^*$ (scalaires \uline{inversibles})}, on a un isom can. d'alg. quadratiques

\[
\mathcal A_Q\simeq\mathcal A_{\lambda Q}
\]

compatible avec les actions sur $E$. Bien sur, il n'en est pas ainsi si on cesse de supposer $\lambda$ inversible (prendre $\lambda=\deleted{\unsure{1}}0$ !).

$3^\circ$) Conditions équivalentes :

a) $\mathcal A_Q\longrightarrow\underline{\operatorname{End}}(E)$ univ. inj. (i.e. inj. sur chaque fibre)

b) $Q$ n'est nulle sur aucune fibre réduite.

c) $\Delta\overset{Q}{\longrightarrow}\operatorname{Quad}(E)$ est univ. injectif

d) $\Sigma_Q$ est strictement fidèle.

(C'est essentiellement le lemme 2 plus haut.) Alors la structure d'Alg. sur $\mathcal A_Q$ est simplement celle induite par $\underline{\operatorname{End}}(E)$.

2.5. Cas tot. dégénéré de la forme \uline{nulle} $Q:E\mathbin{\unsure{-|\longrightarrow}}\deleted{\unsure{L}}\Delta'$. Alors

\[
\mathcal A_Q=\underline O\times\Lambda\qquad(\Lambda=\Delta\otimes L)
\]

avec la structure d'algèbre définie par la condition que $\Lambda$ est de \uline{carré nul}, i.e.
% archive sheet number: 21
